\section{Regeneration in a real direction}\label{sec:regeneration}

Following the regeneration approach of Sznitman and
Zerner~\cite{SznitmanZerner1999}, we construct finite path pieces whose
successive translates are independent.  Their terminal times are selected using the entire future, so the usual
stopping-time Markov property does not establish this independence.  An exact
factorization of finite path weights will do so.  Counting ordinary records
then gives a finite mean projected width, even when the projected heights do
not lie in a discrete subgroup of $\RR$.

\subsection{Departure rows and path weights}

A finite nearest-neighbor word is a sequence
$\gamma=(x_0,\ldots,x_m)$ with $x_{j+1}-x_j\in U$.  Its departure set and
annealed weight are
\[
 \operatorname{Dep}(\gamma)=\{x_0,\ldots,x_{m-1}\},\qquad
 W(\gamma)=\EE\!\left[
       \prod_{j=0}^{m-1}\omega(x_j,x_{j+1}-x_j)\right],
\]
where the expectation in this display is over the environment.  Thus
$W(\gamma)=P_{x_0}(X_j=x_j,\ 0\le j\le m)$.  A row visited repeatedly
contributes repeatedly to the same product inside the expectation.  The
terminal site contributes no row unless the word departed from it earlier.
Translation invariance gives $W(\gamma+z)=W(\gamma)$.

\begin{lemma}[Transfer through unused departure rows]
\label{lem:departure-transfer}
Let $\gamma=(x_0,\ldots,x_m)$ be a finite nearest-neighbor word and let
$f(\omega)$ be a nonnegative measurable function of the rows outside
$\operatorname{Dep}(\gamma)$.  Then
\begin{equation}\label{eq:departure-transfer}
 \EE\!\left[
   \prod_{j=0}^{m-1}\omega(x_j,x_{j+1}-x_j)f(\omega)\right]
       =W(\gamma)\EE[f].
\end{equation}
In particular, two prescribed words with disjoint departure sets have the
same joint weight when sampled under independent annealed laws as when
sampled independently conditional on one shared environment.

The identity also applies when $f(\omega)$ is the quenched probability of an
event for another walk stopped on its first arrival in
$\operatorname{Dep}(\gamma)$.  Events that require this other walk to avoid
that set forever are allowed.
\end{lemma}

\begin{proof}
The product in \eqref{eq:departure-transfer} depends only on rows in
$\operatorname{Dep}(\gamma)$, which are independent of all the rows defining
$f$.  This proves the identity.  For two words, conditional independence in
the shared environment gives the product of their quenched weights, and
\eqref{eq:departure-transfer} factors its environment expectation.

A trajectory stopped on arrival in a set does not use the row at that
arrival site.  Its finite stopped-prefix probabilities therefore depend
only on rows outside the set.  The same is true of every measurable event
of the stopped trajectory, by the monotone class theorem.  This includes
events on which the stopping time is infinite.
\end{proof}

The transfer is made for unclassified path weights.  A condition involving
a quenched exit probability can depend on rows on both sides of the
separation and cannot simply be inserted into the independent side of
\eqref{eq:departure-transfer}.  Later we will first apply this identity and
then classify sites in the shared environment.

\subsection{Fresh records and true cuts}

Fix $v\in\RR^d\setminus\{0\}$, normalized by $\|v\|_\infty=1$, and write
$h(x)=x\cdot v$.  Define
\[
 D_v=\{h(X_n)\ge0\text{ for all }n\ge0\},\qquad
 p_v=P_0(D_v),\qquad
 \widehat P_v=P_0(\,\cdot\mid D_v)
\]
when $p_v>0$.  All statements in Section~\ref{sec:regeneration} concern this fixed real
direction.  We write $\widehat E_v$ for expectation under
$\widehat P_v$.  No arithmetic assumption on the coordinates of $v$ is
imposed.

The next lemma follows the repeated-trial approach of Zerner and
Merkl~\cite[Lemma~4]{ZernerMerkl2001}. We give the argument to specify
which transition rows remain unused at each trial.

\begin{lemma}[Finite height supremum]\label{lem:finite-supremum}
For this fixed real direction, almost surely under $P_0$,
\[
 \sup_{n\ge0}h(X_n)<\infty
 \quad\Longrightarrow\quad h(X_n)\longrightarrow-\infty.
\]
In particular, a finite height interval cannot contain the path forever.
The corresponding statements hold from every translated starting site,
and hold quenched for almost every environment, simultaneously for all
such starting sites.
\end{lemma}

\begin{proof}
Fix nonnegative integers $b,k$, a signed coordinate step $a\in U$ with
$h(a)=1$, and an integer $m>b+k$.  It suffices to rule out
\[
 E_{b,k}=\{h(X_n)\le b\text{ for every }n,
                  \ h(X_n)\ge-k\text{ infinitely often}\}.
\]
Partition the lattice into columns, the cosets of $\ZZ a$.  Every column
has only finitely many sites with height in $[-k,b]$: its successive
heights differ by exactly one, regardless of the other coordinates of
$v$.

First suppose that only finitely many columns are visited at heights at
least $-k$.  On $E_{b,k}$, some fixed site $x$ of height in $[-k,b]$ is
then visited infinitely often.  Fix an environment whose transition
entries are all positive.  At every visit to this particular site the
quenched probability of the script consisting of $m$ steps $a$ is
\[
 q_x(\omega)=\prod_{j=0}^{m-1}\omega(x+ja,a)>0.
\]
Choose the first visit as a trial start, and subsequently choose the
first visit at least $m$ steps after the preceding start.  These are
stopping times.  By the quenched strong Markov property, the probability
that $r$ such trials are finite and all fail is at most
$(1-q_x(\omega))^r$.  Infinitely many visits provide infinitely many
trials, while a successful script would cross above $b$.  Thus the
specified event has quenched probability zero.  A countable union over
$x$ and integration over the environment handle this first case.  There
is no need to bound $q_x(\omega)$ uniformly in $x$ or $\omega$.

For the other case, consider the first visit, for each column, to its
portion of height at least $-k$.  These eligibility decisions are
determined by the observed path prefix.  Select the first eligible time
as a trial start, and thereafter the first eligible time at least $m$
steps after the preceding trial start.  If infinitely many columns have
such visits, infinitely many trial starts remain after these finite
exclusions.  At an eligible time, every earlier visit in the same
column had height below $-k$.  Consequently the $m$ scripted departure
sites, all at height at least $-k$, are distinct and have not appeared
earlier.  Conditional on any fixed eligible prefix, their rows retain
their original independent laws.  The annealed conditional probability
of the script is therefore exactly
\[
 q_a^m,\qquad q_a:=\EE[\omega(0,a)]>0.
\]
The probability that $r$ selected trials are finite and all fail is at
most $(1-q_a^m)^r$.  On $E_{b,k}$ every one fails, so the second case is
also null.  This argument uses no conditioning on the future event
$E_{b,k}$ when evaluating a trial probability.

Taking the countable union over $b,k$ proves the implication: a path
from zero with finite height supremum that does not tend to minus
infinity belongs to some $E_{b,k}$.  Translation gives the same
annealed conclusion from every lattice site.  Integrating the quenched
probabilities of these null events, and taking a countable intersection
over starting sites and integer bounds, gives the last assertion.
\end{proof}

\begin{lemma}\label{lem:positive-nonbacktracking}
If $P_0(A_v)>0$, then $p_v>0$.
\end{lemma}

\begin{proof}
On $A_v$, the sequence $h(X_n)$ tends to $+\infty$ and therefore attains a
global minimum.  Thus $A_v$ is contained in the countable union, over
$n\ge0$ and $x\in\ZZ^d$, of the events
\[
 \{X_n=x,\ h(X_{n+k})\ge h(x)\text{ for every }k\ge0\}.
\]
At least one such event has positive probability.  Put
$q_v(x,\omega)=P_{x,\omega}(h(X_k)\ge h(x)\text{ for all }k\ge0)$.
The quenched Markov property at the deterministic time $n$ gives
\[
 \EE\big[P_{0,\omega}(X_n=x)q_v(x,\omega)\big]>0.
\]
Consequently $\EE[q_v(x,\omega)]>0$, and translation invariance identifies
this expectation with $p_v$.  This argument does not use a Markov property
at the time of the global minimum.
\end{proof}

The ordinary strict-record times are the stopping times
\[
 \rho_0=0,\qquad
 \rho_{i+1}=\inf\{n>\rho_i:h(X_n)>h(X_{\rho_i})\},
\]
with all subsequent times set to infinity if one is infinite.  A finite
$\rho_i$, $i\ge1$, is a \emph{true cut} if
\[
 h(X_n)\ge h(X_{\rho_i})\quad\text{for every }n\ge\rho_i.
\]
Truth of a cut is not measurable at its arrival time.
Whenever true cuts exist, $\tau_1,\tau_2,\ldots$ enumerate them in time
order.

At an ordinary strict record, every earlier departure has height strictly
below the current height.  The closed upper half-space is therefore
unused.  Applying Lemma~\ref{lem:departure-transfer} to each finite record
prefix and translating its endpoint gives
\begin{equation}\label{eq:fresh-record}
 P_0\big(h(X_{\rho_i+k})\ge h(X_{\rho_i})\text{ for all }k\ge0
          \mid\cF_{\rho_i}\big)=p_v
       \quad\text{on }\{\rho_i<\infty\}.
\end{equation}
This is an identity in the annealed path filtration.  Conditional on the
whole environment, the corresponding probability would instead be
$q_v(X_{\rho_i},\omega)$.

For a word $\gamma=(0=x_0,\ldots,x_m=z)$, say that it is a
\emph{regeneration word in direction $v$} if $m\ge1$ and the following two
conditions hold:
\begin{enumerate}
 \item $0\le h(x_j)<h(z)$ for every $0\le j<m$;
 \item whenever $1\le j<m$ is a strict record within the word, there is a
       $k$ with $j<k<m$ and $h(x_k)<h(x_j)$.
\end{enumerate}
The second condition says that every intermediate record has already
failed before the final arrival.  Define the width and radius of such a
word by
\[
 L(\gamma)=h(z)>0,\qquad
 R(\gamma)=\max_{0\le j\le m}\|x_j\|_2.
\]
Both are finite, but the width need not be an integer.

\begin{lemma}[Regeneration words]\label{lem:slabs}
Suppose $p_v>0$.  Under $\widehat P_v$, there are almost surely infinitely
many true cuts
$0=\tau_0<\tau_1<\tau_2<\cdots$, with time zero included as an initial
boundary.  The relative words between consecutive cuts are independent
and identically distributed, with probability law $\nu_v$ given by
\[
 \nu_v(\gamma)=W(\gamma)
\]
on the regeneration words in direction $v$.  The departure heights within
each word belong to $[0,L)$ relative to its starting height.

Under the unconditioned law, on $\{\sup_n h(X_n)=\infty\}$ a first true
strict record exists almost surely.  Conditional on its finite prefix,
the translated future has law $\widehat P_v$.
\end{lemma}

\begin{proof}
We first establish existence using stopping-time trials.  Test the first
ordinary strict record.  If the path later falls below its height, wait
until the first strict record above the entire path seen at that failure
time, and test again.  The candidates and the times at which failures are
detected are stopping times.  By \eqref{eq:fresh-record}, the probability
that the first $k$ tested candidates are finite and fail is at most
$(1-p_v)^k$.  On $\{\sup_n h(X_n)=\infty\}$, a new candidate is available
after every detected failure.  Hence a true cut exists there almost surely.

Under $\widehat P_v$, the height supremum is infinite almost surely by
Lemma~\ref{lem:finite-supremum}: its alternative conclusion of
escape to minus infinity is incompatible with $D_v$.  In particular the
first true cut is finite under $\widehat P_v$.  This application of
Lemma~\ref{lem:finite-supremum} is
valid before any independence or integrability of slabs has been proved.

We now compute its law.  Fix a regeneration word
$\gamma=(0,\ldots,z)$ of length $m$ and a measurable event $B$ for the
translated infinite suffix.  The first true-cut word is $\gamma$ exactly
when this prefix occurs and the suffix from $z$ belongs to $D_v$.
For necessity, an earlier record not invalidated by time $m$ could never
be invalidated by a future remaining above $h(z)$, so it would be an
earlier true cut.  For sufficiency, the word's intermediate records have
all failed, while its terminal record is true.

The finite prefix uses only rows below $h(z)$, whereas its suffix
constrained by $D_v$ uses only rows at or above $h(z)$.  Thus
\[
 P_0\big((X_0,\ldots,X_m)=\gamma,\ \tau_1=m,
         \ (X_{m+k}-z)_{k\ge0}\in B\big)
       =W(\gamma)P_0(B\cap D_v).
\]
The event on the left implies the original $D_v$.  Division by $p_v$
therefore gives
\begin{equation}\label{eq:slab-factorization}
 \widehat P_v\big(\Gamma_1=\gamma,
               \ (X_{\tau_1+k}-X_{\tau_1})_{k\ge0}\in B\big)
          =W(\gamma)\widehat P_v(B).
\end{equation}
Here $\Gamma_1$ denotes the first relative word.  Since $\tau_1$ is almost
surely finite, summing over all such words proves that their weights sum
to one.  Equation~\eqref{eq:slab-factorization} gives an independent suffix
with the original conditioned law, and iteration proves the iid
assertion.  The departure-height claim is part of the word definition.

For the unconditioned first true cut, the same word enumeration allows a
prefix that goes below zero.  Its endpoint is still a strict global
record, and each earlier strict record must have failed before that
endpoint.  The fresh-suffix factor is again $P_0(B\cap D_v)$, so its
conditional translated law is $\widehat P_v$.  No Markov property at a
future-dependent cut has been used.
\end{proof}

We call these relative words \emph{slabs}.  The preceding proof gives
their complete law, not just independence of their displacements.  This
will allow us to arrange whole slabs in different spatial positions.

\subsection{Mean width and the two signs}

The record heights can have arbitrarily small positive increments in a
real direction.  The next argument instead renews in the integer index
of an ordinary record.

\begin{lemma}[Finite mean width]\label{lem:mean-width}
Suppose $p_v>0$ and $\|v\|_\infty=1$.  If $S$ is the number of ordinary
strict records after the initial boundary of a slab with law $\nu_v$,
including its terminal record, then
\[
 0<\mathbb E_{\nu_v}L\le\mathbb E_{\nu_v}S\le p_v^{-1}<\infty.
\]
\end{lemma}

\begin{proof}
For $i\ge1$, let
\[
 B_i=\{\rho_i<\infty,\ h(X_n)\ge0\text{ for }0\le n\le\rho_i\}.
\]
Freshness at $\rho_i$ and cancellation of the conditioning probability
give
\begin{equation}\label{eq:record-renewal-mass}
 \widehat P_v(i\text{ is a true-record index})
   =\frac{P_0(B_i,\ \rho_i\text{ is true})}{p_v}
   =P_0(B_i)\ge p_v.
\end{equation}
The last inequality follows because $D_v$ forces an infinite supremum
and hence reaches every ordinary record index before dropping below zero.

At a cut, its height is the maximum of the entire past.  The subsequent
ordinary records are therefore exactly the records of the translated
suffix.  By Lemma~\ref{lem:slabs}, the successive true-record indices are
sums of iid positive integers $S_1,S_2,\ldots$ with the law of $S$.  Let
\[
 N(n)=\#\{k\ge1:S_1+\cdots+S_k\le n\}.
\]
Summing \eqref{eq:record-renewal-mass} yields
$\widehat E_v[N(n)/n]\ge p_v$.
If $\mathbb E_{\nu_v}S=\infty$, the strong law applied to every bounded
truncation $S_j\wedge a$ gives
$(S_1+\cdots+S_k)/k\to\infty$.  Inverting these sums gives
$N(n)/n\to0$ almost surely.  Since $0\le N(n)/n\le1$, dominated
convergence contradicts the preceding lower bound.  Thus the mean of
$S$ is finite.  The ordinary strong law and the same inversion now give
$N(n)/n\to1/\mathbb E_{\nu_v}S$; dominated convergence proves
$\mathbb E_{\nu_v}S\le1/p_v$.

Each ordinary record height gain is at most one: its final step has
height increment at most $\|v\|_\infty=1$, and starts no higher than the
preceding maximum.  Adding the gains within one slab gives $L\le S$.
Finally, $L>0$ almost surely, so its mean is strictly positive.
\end{proof}

Lemma~\ref{lem:mean-width} controls record count and projected width, not the number of
time steps in a slab.

\begin{lemma}[The two directional signs]\label{lem:sign-union}
For every fixed nonzero real direction $v$, if $P_0(A_v)>0$, then
\[
 P_0(A_v\cup A_{-v})=1.
\]
More precisely, up to null sets, infinite height supremum implies $A_v$
and finite height supremum implies $A_{-v}$.
\end{lemma}

\begin{proof}
Normalize $v$ as above.  Lemma~\ref{lem:positive-nonbacktracking} gives
$p_v>0$.  On infinite height supremum, Lemma~\ref{lem:slabs} provides a
first true cut followed by iid slabs.  Their positive widths have finite,
strictly positive mean by Lemma~\ref{lem:mean-width}, so the successive
base heights tend to infinity.  Every later slab stays at or above its
base, and consequently $h(X_n)\to+\infty$.

The finite-supremum implication is exactly
Lemma~\ref{lem:finite-supremum}.  These two alternatives exhaust
all paths and prove the sign union.
\end{proof}

Thus a probability strictly between zero and one for $A_v$ supplies
positive probabilities for both signs.  In Section~\ref{sec:radius}, that
coexistence assumption will give the stronger spatial estimate
$\mathbb E_{\nu_v}R+\mathbb E_{\nu_{-v}}R<\infty$.
